\chapter{Weighted Finite Automata and Rational Power Series}
\label{ch07:top}

Chapter~6 studied finite automata that answer a yes/no question about each string. A \emph{weighted finite automaton} (WFA) replaces that binary verdict by a numerical weight. The main exact theorem of the chapter is simple and powerful:
\begin{quote}
For a WFA over $\mathbb{R}$ or $\mathbb{C}$, the generating function obtained by summing weights by string length is always rational.
\end{quote}
That places WFAs squarely inside the pole-based world of Chapters~1--4.

Two distinctions are important from the start.
\begin{enumerate}
  \item The definition of a WFA makes sense over an arbitrary semiring, but the resolvent formula $(I-zA)^{-1}$ belongs to the real/complex matrix setting.
  \item A \emph{locally stochastic} WFA is not automatically a proper probability distribution on finite strings. Properness is a separate almost-sure-halting condition.
\end{enumerate}
These two distinctions will prevent several common confusions later in the book.

\section{Semirings}

A \emph{semiring} is a set where one can add and multiply, multiplication distributes over addition, and subtraction need not exist. The examples to keep in mind are:
\begin{itemize}
  \item the Boolean semiring, for ordinary accept/reject automata;
  \item the nonnegative reals, for probabilistic or weighted counting models;
  \item the tropical semiring, for shortest-path style problems.
\end{itemize}
General semiring language is useful because it unifies several automata theories at once. For the analytic results of this chapter, however, we will soon specialize to matrices over $\mathbb{R}$ or $\mathbb{C}$.

\section{Weighted Finite Automata}

\begin{definition}
A \emph{weighted finite automaton} over alphabet $\Sigma$ with $n$ states is a triple
\[
  \mathcal{A} = (\alpha, \{A_\sigma\}_{\sigma \in \Sigma}, \beta),
\]
where
\begin{itemize}
  \item $\alpha$ is a $1 \times n$ row vector of initial weights,
  \item $\beta$ is an $n \times 1$ column vector of final weights,
  \item each $A_\sigma$ is an $n \times n$ transition matrix.
\end{itemize}
For a word $w = \sigma_1 \cdots \sigma_k$, its weight is
\[
  f_{\mathcal A}(w) = \alpha A_{\sigma_1} \cdots A_{\sigma_k} \beta.
\]
The empty word has weight
\[
  f_{\mathcal A}(\varepsilon) = \alpha I \beta = \alpha\beta.
\]
\end{definition}

The matrix dimensions are worth noticing:
\[
  (1\times n)(n\times n)\cdots(n\times n)(n\times 1) = (1\times 1),
\]
so the result is a scalar.

The formula is really a sum over hidden state paths. For a two-letter word $\sigma\tau$,
\[
  \alpha A_\sigma A_\tau \beta
  = \sum_{i,j,k} \alpha_i (A_\sigma)_{ij}(A_\tau)_{jk}\beta_k.
\]
So one sums over all choices of start state $i$, intermediate state $j$, and final state $k$. This is exactly the path-sum interpretation of matrix multiplication.

In the Boolean semiring, the same formula says whether there exists an accepting path labeled by the word, so WFAs genuinely extend ordinary finite automata.

\begin{example}[A Fibonacci WFA]
Let $\Sigma=\{a\}$ and define
\[
  \alpha = \begin{pmatrix}1 & 0\end{pmatrix},
  \qquad
  \beta = \begin{pmatrix}0\\1\end{pmatrix},
  \qquad
  A_a = \begin{pmatrix}0 & 1\\1 & 1\end{pmatrix}.
\]
Then
\[
  f_{\mathcal A}(a^k) = \alpha A_a^k \beta.
\]
A standard induction shows that
\[
  A_a^k = \begin{pmatrix}F_{k-1} & F_k\\ F_k & F_{k+1}\end{pmatrix}
  \qquad (k \ge 1),
\]
so
\[
  f_{\mathcal A}(a^k) = F_k.
\]
Thus a two-state WFA already realizes the Fibonacci sequence.
\end{example}

\section{Local Stochasticity Versus Properness}

Now specialize to nonnegative weights. A WFA is \emph{locally stochastic} if
\begin{itemize}
  \item the initial vector $\alpha$ is a probability distribution on states, and
  \item for each state $q$,
  \[
    \sum_{\sigma \in \Sigma} \sum_{q'} (A_\sigma)_{q,q'} + \beta_q = 1.
  \]
\end{itemize}
The intended interpretation is that from state $q$, the automaton either emits a symbol and changes state, or halts via the final weight $\beta_q$.

But this local condition does \emph{not} by itself guarantee a probability distribution on finite strings.

\begin{example}[A non-halting locally stochastic WFA]
Take one state, one letter $a$, and
\[
  \alpha=(1), \qquad A_a=(1), \qquad \beta=(0).
\]
The local stochasticity condition holds, but the machine loops forever with probability $1$ and halts with probability $0$. So every finite string receives probability $0$, and the total mass on $\Sigma^*$ is also $0$.
\end{example}

This is why we need a second notion.

\begin{definition}
A nonnegative WFA is \emph{proper} if the total weight of all finite strings is exactly $1$:
\[
  \sum_{w \in \Sigma^*} f_{\mathcal A}(w) = 1.
\]
Equivalently, the model halts almost surely.
\end{definition}

Chapter~8 will study properness systematically. For the present chapter, the main point is conceptual: local stochasticity is a step toward probability, not the whole story.

\section{The Length Generating Function is Rational}

From now on we work over $\mathbb{R}$ or $\mathbb{C}$. Define the aggregate matrix
\[
  A = \sum_{\sigma \in \Sigma} A_\sigma.
\]
This forgets which symbol was emitted and keeps only the total one-step weight.

The length generating function of the WFA is
\[
  F(z) = \sum_{w \in \Sigma^*} f_{\mathcal A}(w) z^{|w|}.
\]
Group words by length:
\[
  F(z) = \sum_{k \ge 0} z^k \sum_{|w|=k} f_{\mathcal A}(w).
\]
A word of length $k$ is exactly a $k$-tuple of symbols, so
\[
  \sum_{|w|=k} f_{\mathcal A}(w)
  = \sum_{\sigma_1,\dots,\sigma_k \in \Sigma} \alpha A_{\sigma_1}\cdots A_{\sigma_k}\beta.
\]
By distributivity,
\[
  \sum_{\sigma_1,\dots,\sigma_k \in \Sigma} A_{\sigma_1}\cdots A_{\sigma_k}
  = \left(\sum_{\sigma \in \Sigma} A_\sigma\right)^k
  = A^k.
\]
This step uses only distributivity, so it already makes sense over semirings. Hence
\[
  F(z) = \sum_{k \ge 0} z^k \alpha A^k \beta = \alpha \left(\sum_{k \ge 0} (zA)^k\right) \beta.
\]

Now the real/complex specialization matters. As a \emph{formal} matrix identity,
\[
  (I-zA)\sum_{k \ge 0} (zA)^k = I,
\]
so
\[
  \sum_{k \ge 0} (zA)^k = (I-zA)^{-1}
\]
as a matrix power series. Analytically, the same Neumann series converges for
\[
  |z| < 1/\rho(A),
\]
where $\rho(A)$ is the spectral radius of $A$.
So the same algebraic expression serves two roles: formally it is a power-series identity, and inside the disk of convergence it is an honest analytic rational function.

\begin{theorem}[Rationality of WFA length generating functions]
\label{thm:wfa-rational}
For a WFA over $\mathbb{R}$ or $\mathbb{C}$,
\[
  F(z) = \alpha (I-zA)^{-1} \beta,
  \qquad
  A = \sum_{\sigma \in \Sigma} A_\sigma.
\]
In particular, $F(z)$ is a rational function of $z$.
\end{theorem}

Because
\[
  (I-zA)^{-1} = \frac{\operatorname{adj}(I-zA)}{\det(I-zA)},
\]
each entry is rational. The possible poles occur among the reciprocals of the nonzero eigenvalues of $A$, though cancellations can remove some of them after multiplication by $\alpha$ and $\beta$.

\section{Rational Asymptotics and What Must Be Assumed}

Chapter~4 analyzed rational generating functions in terms of their poles. The correct moral for WFAs is therefore this:

\begin{quote}
The asymptotics are controlled by the dominant pole of the scalar rational function $F(z)=\alpha(I-zA)^{-1}\beta$, not merely by the raw eigenvalue multiplicities of $A$.
\end{quote}

A safe statement is the following.

\begin{proposition}
Suppose $F(z)=\alpha(I-zA)^{-1}\beta$ has a unique dominant pole at $z=1/\lambda$ of order $m$. Then
\[
  [z^n]F(z) \sim C\,\lambda^n n^{m-1}
\]
for some constant $C \ne 0$.
\end{proposition}

This formulation avoids several pitfalls at once:
\begin{itemize}
  \item multiple poles on the spectral circle can create oscillations,
  \item the boundary vectors $\alpha$ and $\beta$ can cancel some eigen-directions,
  \item and the relevant order is the actual pole order of the scalar function, not merely the algebraic multiplicity of an eigenvalue of $A$.
\end{itemize}

For proper stochastic WFAs, the coefficients
\[
  p_n = [z^n]F(z) = \sum_{|w|=n} f_{\mathcal A}(w)
\]
form the length distribution. In the simple aperiodic case with a unique simple dominant pole, this total probability decays geometrically with $n$.

\begin{example}[Fibonacci again]
For the Fibonacci WFA,
\[
  A = A_a = \begin{pmatrix}0 & 1\\1 & 1\end{pmatrix},
\]
so
\[
  F(z) = \alpha(I-zA)^{-1}\beta = \frac{z}{1-z-z^2}.
\]
This is the same generating function as in Chapters~1 and~4, and its dominant pole at $z=1/\varphi$ gives
\[
  [z^n]F(z) = F_n \sim \frac{\varphi^n}{\sqrt{5}}.
\]
The WFA viewpoint reaches the same asymptotic through matrix structure.
\end{example}

\section{Recognizable Versus Rational Series}

There is also a noncommutative version of the story. A formal series over words,
\[
  S = \sum_{w \in \Sigma^*} c_w \, w,
\]
can be viewed as assigning a coefficient directly to each word rather than first collapsing everything to a length variable. Such a series is called \emph{recognizable} if it comes from a WFA, and \emph{rational} if it is built from finite series using sum, concatenation, and star.

\begin{theorem}[Kleene--Sch\"utzenberger, used as a black box]
A noncommutative formal series is recognizable if and only if it is rational.
\end{theorem}

The star operation again requires the analogue of a zero constant term, for the same reason as in Chapter~3: one must avoid infinitely many ways to build the same total length. If we now send each word $w$ to the commutative monomial $z^{|w|}$, concatenation of words becomes multiplication of powers of $z$ because lengths add. That is the bridge from noncommutative rationality to the ordinary rational function $F(z)$ above.

\section{What WFAs Can and Cannot Express Exactly}

The rationality theorem gives an exact limitation. An exact WFA model always has a rational length generating function. Therefore it cannot exactly realize a genuine square-root algebraic length law such as the Catalan/Dyck behavior studied in Chapters~3--5.

This is one way to read the probabilistic hierarchy discussed by Icard \cite{Icard2020}: finite-state weighted models sit at a rational level, while context-free probabilistic models can produce algebraic length behavior above that level. The exact details of that hierarchy are not needed here. The main message is simpler: exact WFA models occupy the rational regime.

\section{Why This Matters for Language Models}

A large language model is not literally a WFA. Still, the WFA viewpoint matters in two different ways.

First, it provides a tractable exact model class. When a sequence model really is finite-state weighted, its length generating function is rational and the pole analysis of Part~I applies exactly.

Second, it provides a research program for approximation. If some finite-state surrogate approximates a language model well on finite strings or finite contexts, that may still be useful. But one must be careful: closeness of distributions or coefficients on a finite range does \emph{not} automatically imply closeness of dominant singularities or asymptotic regimes. Chapter~12 returns to this issue and formulates approximation targets more carefully.

A good rule of thumb is:
\begin{itemize}
  \item exact WFA model $\Rightarrow$ exact rational length generating function;
  \item observed square-root algebraic correction $\Rightarrow$ evidence against an exact WFA description;
  \item useful finite-length approximation does not automatically imply the same asymptotic singularity type.
\end{itemize}
We will use that distinction repeatedly later.

\section*{Further Reading}

\begin{itemize}
  \item \textbf{M. Droste, W. Kuich, and H. Vogler} (eds.), \emph{Handbook of Weighted Automata} (2009) \cite{DrosteKuichVogler2009} --- the standard general reference for weighted automata over semirings.
  \item \textbf{J. Berstel and C. Reutenauer}, \emph{Noncommutative Rational Series with Applications} (2011) \cite{BerstelReutenauer2011} --- for recognizable versus rational series and the noncommutative Kleene--Sch\"utzenberger theorem.
  \item \textbf{T. Icard}, ``Calibrating generative models: the probabilistic Chomsky--Sch\"utzenberger hierarchy'' (2020) \cite{Icard2020} --- for one formal version of the rational-versus-algebraic probabilistic hierarchy.
  \item \textbf{P. Flajolet and R. Sedgewick}, \emph{Analytic Combinatorics} (2009) \cite{FlajoletSedgewick2009} --- for the rational-function asymptotics that make the WFA theorem analytically useful.
\end{itemize}
